\documentclass[12pt,a4paper]{article}

% Packages
\usepackage[left=2.5cm,right=2.5cm,top=2.5cm,bottom=2.5cm]{geometry} % Margins
\usepackage{graphicx} % For including images
\usepackage{float} % For positioning figures
\usepackage{amsmath} % Math symbols
\usepackage{hyperref} % Hyperlinks
\usepackage[utf8]{inputenc} % UTF-8 encoding
\usepackage[T1]{fontenc} % Font encoding
\usepackage{caption} % Better figure captions

% Document Information
\title{Rapport Technique : Apprentissage Auto-Supervisé avec SimCLR sur CIFAR-10}
\author{Projet Deep Learning \newline Encadré par : \textbf{Chahiri}}
\date{Mai 2026}

\begin{document}

% Title Page
\maketitle
\newpage

% Table of Contents
\tableofcontents
\newpage

% Section 1: Introduction
\section{Introduction}
\subsection{Contexte}
L’apprentissage auto-supervisé est un paradigme d’apprentissage non-supervisé qui exploite les données sans étiquettes. Il a émergé comme une solution prometteuse pour des tâches complexes en vision par ordinateur.

\subsection{Problématique}
Comment optimiser la qualité des représentations latentes issues de CIFAR-10 en se basant uniquement sur des techniques contrastives tout en assurant de bonnes performances en évaluation supervisée ?

\subsection{Objectifs}
Les principaux objectifs sont les suivants :
\begin{itemize}
    \item Concevoir un modèle d’apprentissage basé sur SimCLR optimisé pour CIFAR-10.
    \item Évaluer les représentations apprises par une approche linéaire supervisée.
\end{itemize}

\newpage

% Section 2: Méthodologie
\section{Méthodologie}

\subsection{Pipeline SimCLR}
SimCLR (Simple Contrastive Learning Representation) repose sur deux composantes principales :
\begin{itemize}
    \item \textbf{Encodage des images :} Utilisation d’un ResNet-18, adapté à des images petites (32x32 pixels pour CIFAR-10).
    \item \textbf{Projection :} Transformation des représentations encodées en vecteurs d’embedding de dimension 128 via un MLP (Multi-Layer Perceptron).
\end{itemize}

\subsection{Augmentations et Transformations}
Le modèle SimCLR utilise des schémas d’augmentation des données variés pour obtenir des données corrélées :
\begin{itemize}
    \item Recadrages aléatoires (Random Cropping).
    \item Flou gaussien.
    \item Jitters de couleurs.
\end{itemize}
Cette stratégie génère deux versions transformées d’une image donnée pour entraîner le modèle.

\subsection{Optimisation}
La perte contrainte utilisée est la \textbf{NT-Xent Loss} (Normalized Temperature-scaled Cross-Entropy Loss), qui cherche à minimiser la distance entre les représentations positives tout en maximisant celle entre négatives.

\newpage

% Section 3: Implémentation
\section{Implémentation}

SimCLR a été implémenté avec \textbf{PyTorch}. Les principaux fichiers du projet sont :
\begin{itemize}
    \item \textbf{models.py} : Définition de l’architecture.
    \item \textbf{train\_simclr.py} : Entraînement avec NT-Xent Loss.
    \item \textbf{eval\_linear.py} : Évaluation supervisée en ajoutant un classifieur linéaire.
\end{itemize}

\begin{figure}[H]
    \centering
    \includegraphics[width=0.8\textwidth]{/home/Chahiri/Documents/projects/ssl-simclr-cifar10/results/embeddings/umap.png}
    \caption{Visualisation des embeddings via UMAP.}
    \label{fig:umap}
\end{figure}

\newpage

% Section 4: Résultats
\section{Résultats}

\subsection{Embeddings Visualisés}
\textbf{UMAP et t-SNE} sont utilisés pour observer la qualité des représentations produites.

\begin{figure}[H]
    \centering
    \includegraphics[width=0.8\textwidth]{/home/Chahiri/Documents/projects/ssl-simclr-cifar10/results/embeddings/tsne.png}
    \caption{Visualisation des embeddings via t-SNE.}
    \label{fig:tsne}
\end{figure}

\subsection{Courbes de Performance}
Les graphiques ci-dessous présentent les métriques clés obtenues durant le fine-tuning :
\begin{figure}[H]
    \centering
    \includegraphics[width=0.8\textwidth]{/home/Chahiri/Documents/projects/ssl-simclr-cifar10/results/figures/finetune/Train_Loss_step.png}
    \caption{Perte pendant l’entraînement.}
    \label{fig:train_loss}
\end{figure}

\begin{figure}[H]
    \centering
    \includegraphics[width=0.8\textwidth]{/home/Chahiri/Documents/projects/ssl-simclr-cifar10/results/figures/finetune/Test_acc.png}
    \caption{Précision obtenue sur les données de test.}
    \label{fig:test_acc}
\end{figure}

\newpage

% Section 5: Défis et Perspective
\section{Défis et Résolutions}
\begin{itemize}
    \item \textbf{Limitation mémoire GPU :} Réduction de la taille des batchs et ajustements de l’optimiseur.
    \item \textbf{Transformations :} Certaines augmentations (par ex. : flou) ont montré des bénéfices limités.
\end{itemize}

---

% Section 6: Conclusion et Perspectives
\section{Conclusion et Perspectives}

\subsection{Résumé}
Ce projet démontre l’efficacité de SimCLR pour l’aprentissage non-supervisé. Environ \textbf{85\%} de précision est atteinte dans le cadre d’une évaluation supervisée.

\subsection{Futurs Travaux}
\begin{itemize}
    \item Utiliser des modèles plus profonds (ex. : ResNet-50).
    \item Étendre le pipeline à d’autres jeux de données.
\end{itemize}

\end{document}